\originalchapter{屏蔽、声子与电子寿命}
\label{ch:phonon}
\source{8.511第18--21讲。}
\originalsection{密度响应与屏蔽}
外部电势扰动使电子密度改变。为避免符号混淆，定义电子能量扰动$\delta v$及静态极化$\delta n=-\Pi\,\delta v$，$\Pi(q,0)>0$。长波极限$\Pi(0,0)=\partial n/\partial\mu=D(E_F)$。总势为$\delta v_{\rm tot}=\delta v_{\rm ext}+v_q\delta n$，自洽解给
\begin{equation}
\delta v_{\rm tot}=\frac{\delta v_{\rm ext}}{\epsilon(q,\omega)},\qquad \epsilon=1+v_q\Pi,\qquad \epsilon(q,0)=1+\frac{\kappa^2}{q^2},\quad\kappa^2=4\pi e_*^2D(E_F).
\end{equation}
点电荷势因此变为Yukawa形式$e_*^2\e^{-\kappa r}/r$。这里Thomas--Fermi近似忽略了$\Pi$的有限$q$结构，不能产生正确的Friedel长距离振荡。

无相互作用费米子的响应来自一个粒子从已占据态跃到未占据态。以$\chi=-\Pi$的另一常用约定写Lindhard函数，
\begin{equation}
\chi_0(q,\omega)=2\int\frac{\dd^3k}{(2\pi)^3}\frac{f(\epsilon_k)-f(\epsilon_{k+q})}{\hbar\omega+\epsilon_k-\epsilon_{k+q}+\ii0^+}.
\end{equation}
分子是可用初末态的占据差，$\ii0^+$保证推迟因果性。三维零温静态积分给
\begin{equation}
\Pi(q,0)=D(E_F)\left[\frac12+\frac{1-x^2}{4x}\log\left|\frac{1+x}{1-x}\right|\right],\qquad x=\frac q{2k_F}.
\end{equation}
$2k_F$处导数非解析，Fourier变换产生远处的$\cos(2k_Fr)/r^3$密度振荡。在高频长波极限，连续性和惯性给$\epsilon\simeq1-\omega_p^2/\omega^2$，$\omega_p^2=ne^2/(m\epsilon_0)$。零点$\epsilon=0$表示自洽电荷集体振荡，区别于单个粒子--空穴激发。
\originalsection{晶格振动的量子化}
在稳定平衡位置附近，势能为$\frac12\sum_{ij,\alpha\beta}u_{i\alpha}K_{i\alpha,j\beta}u_{j\beta}$。质量加权并Fourier变换得到动力学矩阵$D_{\alpha\beta}(q)$，本征值是$\omega_{q\lambda}^2$。正本征值表示谐振稳定；整体平移要求$q=0$存在零频声学模。每个原胞含$r$原子时有$3r$支，通常3支声学与$3r-3$支光学。

正规坐标量子化为
\begin{equation}
\vct u_j=\frac1{\sqrt N}\sum_{q\lambda}\sqrt{\frac{\hbar}{2M\omega_{q\lambda}}}\vct e_{q\lambda}\left(b_{q\lambda}\e^{\ii q\cdot R_j}+b_{q\lambda}^\dagger\e^{-\ii q\cdot R_j}\right),\quad H_{\rm ph}=\sum_{q\lambda}\hbar\omega_{q\lambda}(b^\dagger b+1/2).
\end{equation}
光学模中的不同质量需放回各基底分量。低温只有$\omega\simeq c_\lambda q$声学模重要，态密度$D_\lambda(\omega)/V=\omega^2/(2\pi^2c_\lambda^3)$。能量积分$\int\hbar\omega D_\lambda(\omega)n_B(\omega)\dd\omega$给
\begin{equation}
\frac{C_V}{V}=\frac{2\pi^2\kb^4T^3}{15\hbar^3}\sum_\lambda c_\lambda^{-3}.
\end{equation}
Debye近似用截止频率使总模数等于$3N$；它是低能声学模型，并不描述所有实际高频支。

离子jellium的裸等离子体频率为$\omega_i^2=n_iZ^2e^2/(M\epsilon_0)$。电子静态屏蔽把它改成$\omega^2(q)=\omega_i^2/\epsilon(q,0)\simeq\omega_i^2q^2/\kappa^2$。离子本来有隙的纵电荷振荡因此在长波成为声波，速度$c_L=\omega_i/\kappa$。这是电子可比离子快地调整密度的绝热近似。
\originalsection{电子--声子作用与幂律}
晶格位移改变电子势，保留一阶得到$H_{e\text{-}ph}=\sum_{kq\sigma}g_qc_{k+q,\sigma}^\dagger c_{k\sigma}(b_q+b_{-q}^\dagger)$。声学形变势常有$|g_q|^2\propto q$，因为应变含$qu$，而量子位移$u\propto\omega_q^{-1/2}$。零温费米能上方$E$的电子可发射声子，Golden rule为
\begin{equation}
\frac1\tau=\frac{2\pi}{\hbar}\int\frac{\dd^3q}{(2\pi)^3}|g_q|^2[1-f(\epsilon_{k-q})]\delta(\epsilon_k-\epsilon_{k-q}-\hbar c_sq).
\end{equation}
能量守恒限制$q\lesssim E/(\hbar c_s)$。在$E\ll\hbar\omega_D$且$v_F\gg c_s$时，角积分的$\delta$函数给$1/(\hbar v_Fq)$，结合相空间$q^2\dd q$和$|g_q|^2\propto q$，剩$\int^{q_{\max}}q^2\dd q\propto E^3$。这给单粒子衰减率的$E^3$幂律。输运还须乘$1-\cos\theta_{kk'}\simeq q_\perp^2/(2k_F^2)$，故动量弛豫率多两个$E$次幂。把热激发尺度取$\kb T$得到低温电阻的$T^5$规律，但数值系数、角向限制与材料参数须由具体$g_q$计算。

声子传播子的静态低频贡献为负：对谐振坐标完成平方或积分掉它，得到$V_{\rm eff}(q,\omega)=|g_q|^2\,2\hbar\omega_q/[(\hbar\omega+\ii0^+)^2-(\hbar\omega_q)^2]$。$|\omega|\ll\omega_q$时这是吸引，频率高过声子尺度时符号与强度改变。其迟滞使短时库仑排斥与低频配对吸引的作用窗口不同。对电子实自能展开$\Sigma'(\omega)\simeq-\lambda\hbar\omega$得到质量增强$1+\lambda$；比热相应增强，而自旋响应还含顶角修正，不能仅以质量替换判断所有观测量。
